\documentclass[10pt]{extarticle}
\usepackage[a4paper,margin=1.5cm]{geometry}
\usepackage[T1]{fontenc}
\usepackage{amsmath,amssymb,mathtools}
\usepackage{graphicx}
\usepackage[export]{adjustbox}
\usepackage{multicol}
\usepackage{enumitem}
\usepackage{booktabs,array,tabularx}
\setlength{\columnsep}{0.7cm}
\setlength{\parindent}{0pt}
\setlength{\parskip}{3pt}
\pagestyle{empty}
\makeatletter
\newcommand{\qfigdef}[4]{\@namedef{qf@file@#1}{#2}\@namedef{qf@w@#1}{#3}\@namedef{qf@wide@#1}{#4}}
\newcommand{\qf@check}[1]{\@ifundefined{qf@file@#1}{\PackageError{score2}{Unknown figure id #1}{}}{}}
\newcommand{\qf@img}[4][]{\includegraphics[width=#2\dimexpr\@nameuse{qf@w@#4}pt\relax,max width=#3,max height=0.45\textheight,#1]{\@nameuse{qf@file@#4}}}
\newif\ifqf@full
\newif\ifqf@boxed
\newsavebox{\qf@box}
\newcommand{\qf@open}{\global\qf@fulltrue\par\global\qf@boxedtrue
  \setbox\qf@box\vbox\bgroup\hsize\textwidth\linewidth\textwidth}
\newcommand{\qf@emit}[1]{\par\egroup\global\qf@boxedfalse
  \ifdim\dimexpr\ht\qf@box+\dp\qf@box\relax<#1\textheight\noindent\usebox{\qf@box}\par\else\unvbox\qf@box\fi}
\newcommand{\qfullsplit}{\ifqf@boxed\qf@emit{0.85}\fi}
\newcommand{\qf@close}{\ifqf@boxed\qf@emit{0.45}\fi\global\qf@fullfalse}
\newcommand{\qfullbegin}{\end{multicols}\qf@open}
\newcommand{\qfullend}{\qf@close\begin{multicols}{2}}
\newcommand{\qfullstart}{\qf@open}
\newcommand{\qfulldone}{\qf@close}
\newcommand{\qfig}[2][1]{\qf@check{#2}%
  \ifqf@full\qf@blk{#1}{#2}\else
  \ifnum\@nameuse{qf@wide@#2}=1 \expandafter\@firstoftwo\else\expandafter\@secondoftwo\fi
  {\end{multicols}\begin{center}\qf@img{#1}{\textwidth}{#2}\end{center}\begin{multicols}{2}}%
  {\qf@blk{#1}{#2}}\fi}
\newcommand{\qf@blk}[2]{\par\smallskip{\centering\qf@img{#1}{\linewidth}{#2}\par}\smallskip}
\newcommand{\qfigi}[2][1]{\qf@check{#2}\qf@img[valign=c]{#1}{\linewidth}{#2}}
\makeatother
\newcommand{\papersec}[1]{\par\bigskip{\centering\Large\bfseries PAPER-#1\par}\medskip}
\newcommand{\partsec}[2]{\par\medskip\noindent{\large\bfseries PART-#1: #2}\par\smallskip}
\newcommand{\sect}[1]{\par\smallskip\noindent\textbf{#1}\par}
\newcommand{\qn}[1]{\par\smallskip\noindent\textbf{#1.}~}
\newcommand{\paraq}[2]{\par\smallskip\noindent\textbf{Paragraph for Questions #1 and #2}\par}
\newenvironment{qoptions}{\begin{itemize}[leftmargin=1.8em,itemsep=1pt,topsep=2pt,label={}]}{\end{itemize}}
\newcommand{\opt}[2]{\item[(#1)] #2}
\newcolumntype{L}[1]{>{\raggedright\arraybackslash}p{#1}}
\newcommand{\fwstart}{\end{multicols}\par\noindent\begin{minipage}{\textwidth}\small\setlength{\tabcolsep}{4pt}\renewcommand{\arraystretch}{1.5}}
\newcommand{\fwend}{\end{minipage}\par\medskip\begin{multicols}{2}}
% figure registry (generated by score2_tools.py finalize)
\qfigdef{p1_phy_s1_q01_1}{SRG (Enthusiast) - Test 12 (19 Sep 2021)-img/p1_phy_s1_q01_1.png}{198.1}{0}
\qfigdef{p1_phy_s1_q04_1}{SRG (Enthusiast) - Test 12 (19 Sep 2021)-img/p1_phy_s1_q04_1.png}{178.5}{0}
\qfigdef{p1_phy_s1_q06_1}{SRG (Enthusiast) - Test 12 (19 Sep 2021)-img/p1_phy_s1_q06_1.png}{152.1}{0}
\qfigdef{p1_phy_s2_q08_1}{SRG (Enthusiast) - Test 12 (19 Sep 2021)-img/p1_phy_s2_q08_1.png}{145.5}{0}
\qfigdef{p1_phy_s2_q09_1}{SRG (Enthusiast) - Test 12 (19 Sep 2021)-img/p1_phy_s2_q09_1.png}{113.9}{0}
\qfigdef{p1_phy_s2_q10_1}{SRG (Enthusiast) - Test 12 (19 Sep 2021)-img/p1_phy_s2_q10_1.png}{125.6}{0}
\qfigdef{p1_phy_s2_q11_1}{SRG (Enthusiast) - Test 12 (19 Sep 2021)-img/p1_phy_s2_q11_1.png}{185.0}{0}
\qfigdef{p1_phy_s2_q12_1}{SRG (Enthusiast) - Test 12 (19 Sep 2021)-img/p1_phy_s2_q12_1.png}{151.8}{0}
\qfigdef{p1_phy_s3_q13_1}{SRG (Enthusiast) - Test 12 (19 Sep 2021)-img/p1_phy_s3_q13_1.png}{148.5}{0}
\qfigdef{p1_phy_s3_q14_1}{SRG (Enthusiast) - Test 12 (19 Sep 2021)-img/p1_phy_s3_q14_1.png}{93.8}{0}
\qfigdef{p1_phy_s3_q15_1}{SRG (Enthusiast) - Test 12 (19 Sep 2021)-img/p1_phy_s3_q15_1.png}{156.5}{0}
\qfigdef{p1_phy_s3_q16_1}{SRG (Enthusiast) - Test 12 (19 Sep 2021)-img/p1_phy_s3_q16_1.png}{93.8}{0}
\qfigdef{p1_phy_s3_q17_1}{SRG (Enthusiast) - Test 12 (19 Sep 2021)-img/p1_phy_s3_q17_1.png}{166.0}{0}
\qfigdef{p1_chem_s1_q03_A}{SRG (Enthusiast) - Test 12 (19 Sep 2021)-img/p1_chem_s1_q03_A.png}{81.4}{0}
\qfigdef{p1_chem_s1_q03_2}{SRG (Enthusiast) - Test 12 (19 Sep 2021)-img/p1_chem_s1_q03_2.png}{111.5}{0}
\qfigdef{p1_chem_s1_q03_C}{SRG (Enthusiast) - Test 12 (19 Sep 2021)-img/p1_chem_s1_q03_C.png}{75.9}{0}
\qfigdef{p1_chem_s1_q03_3}{SRG (Enthusiast) - Test 12 (19 Sep 2021)-img/p1_chem_s1_q03_3.png}{180.0}{0}
\qfigdef{p1_chem_s1_q03_4}{SRG (Enthusiast) - Test 12 (19 Sep 2021)-img/p1_chem_s1_q03_4.png}{156.6}{0}
\qfigdef{p1_chem_s1_q04_A}{SRG (Enthusiast) - Test 12 (19 Sep 2021)-img/p1_chem_s1_q04_A.png}{96.9}{0}
\qfigdef{p1_chem_s1_q04_B}{SRG (Enthusiast) - Test 12 (19 Sep 2021)-img/p1_chem_s1_q04_B.png}{96.9}{0}
\qfigdef{p1_chem_s1_q04_C}{SRG (Enthusiast) - Test 12 (19 Sep 2021)-img/p1_chem_s1_q04_C.png}{99.5}{0}
\qfigdef{p1_chem_s1_q04_D}{SRG (Enthusiast) - Test 12 (19 Sep 2021)-img/p1_chem_s1_q04_D.png}{109.7}{0}
\qfigdef{p1_chem_s1_q05_1}{SRG (Enthusiast) - Test 12 (19 Sep 2021)-img/p1_chem_s1_q05_1.png}{86.8}{0}
\qfigdef{p1_chem_s1_q06_1}{SRG (Enthusiast) - Test 12 (19 Sep 2021)-img/p1_chem_s1_q06_1.png}{199.6}{0}
\qfigdef{p1_chem_s1_q06_A}{SRG (Enthusiast) - Test 12 (19 Sep 2021)-img/p1_chem_s1_q06_A.png}{92.1}{0}
\qfigdef{p1_chem_s1_q06_B}{SRG (Enthusiast) - Test 12 (19 Sep 2021)-img/p1_chem_s1_q06_B.png}{92.3}{0}
\qfigdef{p1_chem_s1_q06_C}{SRG (Enthusiast) - Test 12 (19 Sep 2021)-img/p1_chem_s1_q06_C.png}{97.1}{0}
\qfigdef{p1_chem_s1_q06_D}{SRG (Enthusiast) - Test 12 (19 Sep 2021)-img/p1_chem_s1_q06_D.png}{83.8}{0}
\qfigdef{p1_chem_s2_q10_1}{SRG (Enthusiast) - Test 12 (19 Sep 2021)-img/p1_chem_s2_q10_1.png}{234.1}{0}
\qfigdef{p1_chem_s2_q10_A}{SRG (Enthusiast) - Test 12 (19 Sep 2021)-img/p1_chem_s2_q10_A.png}{98.6}{0}
\qfigdef{p1_chem_s2_q10_B}{SRG (Enthusiast) - Test 12 (19 Sep 2021)-img/p1_chem_s2_q10_B.png}{106.5}{0}
\qfigdef{p1_chem_s2_q11_1}{SRG (Enthusiast) - Test 12 (19 Sep 2021)-img/p1_chem_s2_q11_1.png}{123.9}{0}
\qfigdef{p1_chem_s2_q12_A}{SRG (Enthusiast) - Test 12 (19 Sep 2021)-img/p1_chem_s2_q12_A.png}{203.8}{0}
\qfigdef{p1_chem_s2_q12_B}{SRG (Enthusiast) - Test 12 (19 Sep 2021)-img/p1_chem_s2_q12_B.png}{58.0}{0}
\qfigdef{p1_chem_s2_q12_C}{SRG (Enthusiast) - Test 12 (19 Sep 2021)-img/p1_chem_s2_q12_C.png}{58.9}{0}
\qfigdef{p1_chem_s2_q12_D}{SRG (Enthusiast) - Test 12 (19 Sep 2021)-img/p1_chem_s2_q12_D.png}{58.9}{0}
\qfigdef{p1_chem_s3_q16_A}{SRG (Enthusiast) - Test 12 (19 Sep 2021)-img/p1_chem_s3_q16_A.png}{369.4}{1}
\qfigdef{p1_chem_s3_q16_2}{SRG (Enthusiast) - Test 12 (19 Sep 2021)-img/p1_chem_s3_q16_2.png}{392.5}{1}
\qfigdef{p1_chem_s3_q17_1}{SRG (Enthusiast) - Test 12 (19 Sep 2021)-img/p1_chem_s3_q17_1.png}{128.1}{0}
\qfigdef{p1_chem_s3_q17_2}{SRG (Enthusiast) - Test 12 (19 Sep 2021)-img/p1_chem_s3_q17_2.png}{56.9}{0}
\qfigdef{p1_chem_s3_q17_3}{SRG (Enthusiast) - Test 12 (19 Sep 2021)-img/p1_chem_s3_q17_3.png}{54.4}{0}
\qfigdef{p1_chem_s3_q17_4}{SRG (Enthusiast) - Test 12 (19 Sep 2021)-img/p1_chem_s3_q17_4.png}{58.3}{0}
\raggedcolumns\makeatletter\@beginparpenalty=10000\makeatother

\begin{document}
\begin{center}{\Large\bfseries SRG (Enthusiast) --- Test 12 (19 Sep 2021)}\end{center}
\begin{multicols}{2}
\partsec{1}{PHYSICS}
\sect{SECTION-I (i)}
\qn{1} Find the magnetic induction vector at origin $O$ due to the current carrying wire configuration as shown.
\qfig{p1_phy_s1_q01_1}
\begin{qoptions}
\opt{A}{$\dfrac{\mu_0 I}{4\pi a}\left[\hat i-\dfrac{\pi}{2}\hat j-2\hat k\right]$}
\opt{B}{$\dfrac{\mu_0 I}{2\pi a}\left[\hat i-\dfrac{\pi}{2}\hat j-2\hat k\right]$}
\opt{C}{$\dfrac{\mu_0 I}{4\pi a}\left[-\hat i+\dfrac{\pi}{2}\hat j+2\hat k\right]$}
\opt{D}{$\dfrac{\mu_0 I}{2\pi a}\left[\hat i+\dfrac{\pi}{2}\hat j+2\hat k\right]$}
\end{qoptions}
\qn{2} A capacitor of capacitance $C$ is connected to a battery of EMF $E$ for a long time and then disconnected. The charged capacitor is then connected across a long solenoid having $n$ turns per meter in its closely packed winding on its core. After connections it is found that the voltage across the capacitor drops to $E/\eta$ in a time $\Delta t$. In this period estimate the average magnetic induction at the centre of solenoid.
\begin{qoptions}
\opt{A}{$\dfrac{\mu_0 nCE}{\Delta t}\left(1-\dfrac1\eta\right)$}
\opt{B}{$\dfrac{\mu_0 nCE}{\Delta t\,\eta}$}
\opt{C}{$\dfrac{\mu_0 CE}{n\Delta t}\left(1-\dfrac1\eta\right)$}
\opt{D}{$\dfrac{CE}{\mu_0 n\Delta t}\left(1-\dfrac1\eta\right)$}
\end{qoptions}
\qn{3} An electron with velocity $v$ enters a long solenoid of radius $R$ having magnetic field $B$ radially. Find the time spent by the electron inside the solenoid.
\begin{qoptions}
\opt{A}{$\dfrac{2m}{eB}\tan^{-1}\!\left(\dfrac{eBR}{mv}\right)$}
\opt{B}{$\dfrac{m}{eB}\tan^{-1}\!\left(\dfrac{eBR}{mv}\right)$}
\opt{C}{$\dfrac{eB}{m}\sin^{-1}\!\left(\dfrac{eBR}{mv}\right)$}
\opt{D}{$\dfrac{2m}{eB}\sin^{-1}\!\left(\dfrac{eBR}{m}\right)$}
\end{qoptions}
\qn{4} Two adiabatic containers contain $n$ moles of He at temperature $T_1$ and $n$ moles of H$_2$ at temperature $T_2$ ($T_2<T_1$). They are joined by a conducting rod of length $L$, area $A$ and conductivity $k$. Find the time after which the temperature difference becomes half of the initial.
\qfig{p1_phy_s1_q04_1}
\begin{qoptions}
\opt{A}{$\dfrac{5nRL}{8KA}\,\ell n3$}
\opt{B}{$\dfrac{15}{8}\dfrac{RL}{nKA}\,\ell n4$}
\opt{C}{$\dfrac{15nRL}{16KA}\,\ell n2$}
\opt{D}{$\dfrac58\dfrac{nRL}{nKA}\,\ell n2$}
\end{qoptions}
\qn{5} A liquid $A$ of mass $m$ at $90^\circ$C cools to $30^\circ$C in 5 min in a room at $20^\circ$C. Another liquid $B$ starting at $50^\circ$C takes 5 min to reach $30^\circ$C. If the two liquids are mixed at $90^\circ$C and $50^\circ$C respectively, find the range of final temperature (Newton's law of cooling).
\begin{qoptions}
\opt{A}{$70^\circ$C}
\opt{B}{$60^\circ$C}
\opt{C}{$54^\circ$C}
\opt{D}{$66^\circ$C}
\end{qoptions}
\qn{6} A solid sphere rolls, its centre of mass $C$ being at a distance $d$ from the centre $O$, with angular velocity $\omega$, radius $R$ and mass $m$. Find the angular acceleration.
\qfig{p1_phy_s1_q06_1}
\begin{qoptions}
\opt{A}{$\dfrac52\dfrac{gd}{R^2}$}
\opt{B}{$\dfrac57\dfrac{gd}{R^2}$}
\opt{C}{$\dfrac55\dfrac{d(g+\omega^3R)}{R^2}$}
\opt{D}{$\dfrac57\dfrac{d(g+\omega^3R)}{R^2}$}
\end{qoptions}
\sect{SECTION-I (ii)}
\qn{7} A particle of specific charge $\alpha$ and mass 1 unit moves with velocity $\vec V=V_0(\hat i+\hat k)$ in a magnetic field $\vec B=-B_0\hat k$. Choose the correct option(s).
\begin{qoptions}
\opt{A}{At $t=\dfrac{\pi}{\alpha B_0}$ the velocity of the particle is $-V_0(\hat i-\hat k)$}
\opt{B}{At $t=\dfrac{\pi}{4\alpha B_0}$ the speed of the particle is $V_0$}
\opt{C}{At $t=\dfrac{2\pi}{\alpha B_0}$ the magnitude of displacement is more than $\dfrac{2V_0}{\alpha B_0}$}
\opt{D}{At $t=\dfrac{2\pi}{\alpha B_0}$ the distance travelled is less than $\dfrac{2\sqrt2\,\pi V_0}{\alpha B_0}$}
\end{qoptions}
\qn{8} The figure contains an infinite slab having current per unit area of $\vec J_1=J\hat k$ between the infinite planes at $x=-b$ and $x=b$. Slightly to the right of $x=b$, an infinite thin sheet is kept. It carries a current per unit length $\vec k_2=2bJ(-\hat k)$. Choose the correct option(s).
\qfig{p1_phy_s2_q08_1}
\begin{qoptions}
\opt{A}{The field due to sheet at any point is $\mu_0 Jb$}
\opt{B}{The field due to sheet at any point is $\dfrac{\mu_0 Jb}{2}$}
\opt{C}{The magnetic field due to slab at a point inside the slab must be independent of $y$ co-ordinate and $y$ component of field must be a function of $x$}
\opt{D}{Magnetic field at a point $(x,y)$ inside slab, due to slab is $\mu_0 Jx$}
\end{qoptions}
\qn{9} A charged particle having its charge to mass ratio as $\beta$ goes in a conical pendulum of length $L$ making an angle $\theta$ with vertical and angular velocity $\omega$. If a magnetic field $B$ is directed vertically downwards:
\qfig{p1_phy_s2_q09_1}
\begin{qoptions}
\opt{A}{$B=\dfrac1\beta\left[\omega-\dfrac{g}{\omega L\cos\theta}\right]$}
\opt{B}{Angular momentum of the particle about the point of suspension remains constant}
\opt{C}{If the direction of $B$ were reversed maintaining same $\omega$ and $L$, $\theta$ will remain unchanged}
\opt{D}{Rate of change of angular momentum of the particle about the point of suspension is not a constant vector.}
\end{qoptions}
\qn{10} A mass spectrometer is a device which selects particles of equal mass. An ion with an electric charge $q>0$ starts at rest from a source $S$ and is accelerated through a potential difference $V$. It passes through a hole into a region of constant magnetic field $\vec B$ perpendicular to the plane of the paper as shown in the figure. The particle is deflected by the magnetic field and emerges through the bottom hole at a distance $d$ from the top hole. The mass of the particle is:
\qfig{p1_phy_s2_q10_1}
\begin{qoptions}
\opt{A}{$\dfrac{qBd}{V}$}
\opt{B}{$\dfrac{qB^2d^2}{4V}$}
\opt{C}{$\dfrac{qB^2d^2}{8V}$}
\opt{D}{$\dfrac{qBd}{2V}$}
\end{qoptions}
\qn{11} A current carrying conductor is in the form of a sine curve as shown, which carries current $I$. If the equation of this curve is $Y=2\sin\!\left(\dfrac{\pi x}{L}\right)$ and a uniform magnetic field $B$ exists in space:
\qfig{p1_phy_s2_q11_1}
\begin{qoptions}
\opt{A}{Force on wire is $BIL$ if field is along $Z$-axis}
\opt{B}{Force on wire is $2BIL$ if field is along $Z$-axis}
\opt{C}{Force on wire is zero if field is along $X$-axis}
\opt{D}{Force on wire is $BIL$ if field is in the $XZ$ plane making an angle $30^\circ$ with $X$-axis}
\end{qoptions}
\qn{12} In the circuit shown in the figure, the switch `S' is closed at time $t=0$. The current in branch $AB$ is represented by $z$ and is taken to be positive when it is from $A$ to $B$.
\begin{qoptions}
\opt{A}{The value of $z$ immediately after the switch is closed is $\dfrac VR$}
\opt{B}{The value of $z$ infinite time after the switch is closed is $\dfrac V{3R}$}
\opt{C}{$z$ as a function of time $t$ is $z=\dfrac{4V}{3R}\left(e^{\frac{-3t}{CR}}-\dfrac14\right)$}
\opt{D}{At a time $t_0=\dfrac{2CR}{3}\,\ell n(2)$ the current $z$ becomes zero}
\end{qoptions}
\qfig{p1_phy_s2_q12_1}
\sect{SECTION-II}
\qn{13} All capacitors in the network given below are identical with capacitance of each being $1\,\mu$C. Find the charge on the capacitor marked as $C$ (in $\mu$C).
\qfig{p1_phy_s3_q13_1}
\qn{14} In the given circuit it is known that the capacitor $A$ has a capacitance of $2\,\mu$F and carries a charge of $40\,\mu$C. Capacitor $C$ has a capacitance of $6\,\mu$F and carries a charge of $180\,\mu$C. The positive plate of both capacitors has been indicated in the figure. Capacitance of capacitor $B$ is $3\,\mu$F. Calculate charge on $B$ (in $\mu$C) after the switch $S$ is closed.
\qfig{p1_phy_s3_q14_1}
\qn{15} A wooden disc of mass $M$ and radius $R$ has a single loop of wire wound on its circumference. It is mounted on a massless rod of length $d$. The ends of the rod are supported at its ends so that the rod is horizontal and disc is vertical. A uniform magnetic field $B_0$ exists in vertically upward direction. When a current $I$ is given to the wire one end of the rod leaves the support. If least value of $I=\dfrac{Mg\,d}{N\pi R^2 B_0}$, find $N$.
\qfig{p1_phy_s3_q15_1}
\qn{16} A uniform ring of mass $M$ and radius $R$ carries a current $I$ (see figure). The ring is suspended using two identical strings $OA$ and $OB$. There exists a uniform horizontal magnetic field $B_0$ parallel to the diameter $AB$ of the ring. Calculate ratio of tension $\left(\dfrac{T_{OA}}{T_{OB}}\right)$ in the two strings. [Given $\theta=60^\circ$]

Given $\left[\pi IRB_0=\dfrac{Mg}{2}\right]$
\qfig{p1_phy_s3_q16_1}
\qn{17} The container $A$ contains ice at $0^\circ$C. A conducting uniform rod $PQ$ of length $4R$ is used to transfer heat to the ice in the container. The end $P$ of the rod is maintained at $100^\circ$C and the other end $Q$ is kept inside container $A$. The complete ice melts in 23 minutes. In another experiment, two conductors in shape of quarter circle of radii $2R$ and $R$ are welded to the conductor $PQ$ at $M$ and $N$ respectively and their other ends are inserted inside the container $A$. All conductors are made of same material and have same cross sectional area. Once again the end $P$ is maintained at $100^\circ$C and this time the complete ice melts in $t$ minute. Find $t$. Assume no heat loss from the curved surface of the rods. [Take $\pi\simeq3.0$]
\qfig{p1_phy_s3_q17_1}
\qn{18} A spherical shell is kept in an atmosphere at temperature $T_0$. The wavelength corresponding to maximum intensity of radiation for the shell is $\lambda_0$. A point source of constant power is switched on inside the shell. The power radiated by the source is $P=0.4\sigma SeT_0^4$ where $S$, $e$ and $\sigma$ are outer surface area of the shell, emissivity of the outer surface of the shell and Stefan's constant respectively. Calculate the new wavelength $(\lambda)$ corresponding to the maximum intensity of radiation from the shell. Assume that change in temperature $(\Delta T)$ of the shell is small compared to the ambient temperature $T_0$. If $\lambda=\dfrac{2N}{11}\lambda_0$, find `$N$'.
\partsec{2}{CHEMISTRY}
\sect{SECTION-I (i)}
\qn{1} The dissociation constant of $\mathrm{NH_3}$ at $27^\circ$C from the following data:

$\mathrm{NH_3+H^+\rightleftharpoons NH_4^+}$; $\Delta H^\circ=-52.21$ kJ/mol; $\Delta S^\circ=+1.6$ JK$^{-1}$

$\mathrm{H_2O\rightleftharpoons H^++OH^-}$; $\Delta H^\circ=54.70$ kJ/mol; $\Delta S^\circ=-76.3$ JK$^{-1}$ mol$^{-1}$
\begin{qoptions}
\opt{A}{$e^{10}$}
\opt{B}{$e^{-10}$}
\opt{C}{$e^{-8}$}
\opt{D}{$e^{-9}$}
\end{qoptions}
\qn{2} Calculate the formation constant for the reaction of a tripositive metal ion with thiocyanate ions to form the mono complex if the initial metal concentration in the solution is $2\times10^{-3}$ M, the initial $\mathrm{SCN^-}$ concentration is $1.51\times10^{-3}$ M and the free $\mathrm{SCN^-}$ concentration is $1.0\times10^{-5}$ M.
\begin{qoptions}
\opt{A}{$7.55\times10^{4}$}
\opt{B}{$3\times10^{5}$}
\opt{C}{$3.33\times10^{-6}$}
\opt{D}{$1.5\times10^{5}$}
\end{qoptions}
\qn{3} Among the below given reactions, backward reaction is highly favoured in :
\begin{qoptions}
\opt{A}{\qfigi{p1_chem_s1_q03_A}}
\opt{B}{\qfigi{p1_chem_s1_q03_2}}
\opt{C}{\qfigi{p1_chem_s1_q03_C}}
\opt{D}{\qfigi{p1_chem_s1_q03_3}}
\end{qoptions}
\qn{4} \qfigi{p1_chem_s1_q03_4}

The more possible product (90\% product) of the above reaction is
\begin{qoptions}
\opt{A}{\qfigi{p1_chem_s1_q04_A}}
\opt{B}{\qfigi{p1_chem_s1_q04_B}}
\opt{C}{\qfigi{p1_chem_s1_q04_C}}
\opt{D}{\qfigi{p1_chem_s1_q04_D}}
\end{qoptions}
\qn{5} A tertiary alcohol \textbf{H} upon acid catalysed dehydration gives a product \textbf{I}. Ozonolysis of \textbf{I} leads to compounds \textbf{J} and \textbf{K}. Compound \textbf{J} upon reaction with KOH gives benzyl alcohol and compound \textbf{L}, whereas \textbf{K} on reaction with KOH gives only \textbf{M}.

\qfig{p1_chem_s1_q05_1}

The structure of compounds J, K and L respectively, are
\begin{qoptions}
\opt{A}{$\mathrm{PhCOCH_3,\ PhCH_2COCH_3}$ and $\mathrm{PhCH_2COO^-K^+}$}
\opt{B}{$\mathrm{PhCHO,\ PhCH_2CHO}$ and $\mathrm{PhCOO^-K^+}$}
\opt{C}{$\mathrm{PhCOCH_3,\ PhCH_2CHO}$ and $\mathrm{CH_3COO^-K^+}$}
\opt{D}{$\mathrm{PhCHO,\ PhCOCH_3}$ and $\mathrm{PhCOO^-K^+}$}
\end{qoptions}
\qn{6} \qfigi{p1_chem_s1_q06_1}
\begin{qoptions}
\opt{A}{\qfigi{p1_chem_s1_q06_A}}
\opt{B}{\qfigi{p1_chem_s1_q06_B}}
\opt{C}{\qfigi{p1_chem_s1_q06_C}}
\opt{D}{\qfigi{p1_chem_s1_q06_D}}
\end{qoptions}
\sect{SECTION-I (ii)}
\qn{7} Which of the following statement(s) is/are correct?
\begin{qoptions}
\opt{A}{The pH of $10^{-8}$ M $\mathrm{HCl}$ solution is 8}
\opt{B}{The conjugate base of $\mathrm{H_2PO_4^-}$ is $\mathrm{HPO_4^{2-}}$}
\opt{C}{Autoprotolysis constant of water increases with the increase in temperature.}
\opt{D}{When a solution of a weak monoprotic acid is titrated against a strong base, at half neutralization point, $\mathrm{pH}=\mathrm{p}K_a/2$}
\end{qoptions}
\qn{8} $\mathrm{H_2CO_3}$ ionizes as follows

$\mathrm{H_2CO_3\rightleftharpoons H^++HCO_3^-}$; $K_1=4.0\times10^{-6}$

$\mathrm{HCO_3^-\rightleftharpoons H^++CO_3^{2-}}$; $K_2=5\times10^{-11}$

Which of the following information is correctly related to 0.5 M $\mathrm{Na_2CO_3}$ solution ?
\begin{qoptions}
\opt{A}{The degree of hydrolysis of $\mathrm{Na_2CO_3}$ is 0.02}
\opt{B}{The pH of solution is 2.0.}
\opt{C}{The pOH of solution is 2.0.}
\opt{D}{$[\mathrm{H_2CO_3}]=2.5\times10^{-9}$ M}
\end{qoptions}
\qn{9} A certain indicator (an organic dye) has $\mathrm{p}K_a=5$. For which of the following titrations may it be suitable ?
\begin{qoptions}
\opt{A}{Acetic acid against NaOH.}
\opt{B}{Aniline hydrochloride against NaOH}
\opt{C}{Sodium acetate against HCl.}
\opt{D}{Barium hydroxide against oxalic acid}
\end{qoptions}
\qn{10} \qfigi[0.9]{p1_chem_s2_q10_1}

The major product is
\begin{qoptions}
\opt{A}{\qfigi{p1_chem_s2_q10_A}}
\opt{B}{\qfigi{p1_chem_s2_q10_B}}
\opt{C}{Diastereomers will be formed}
\opt{D}{Racemic mixture}
\end{qoptions}
\qn{11} The correct statement regarding Cannizaro reaction is
\begin{qoptions}
\opt{A}{$\mathrm{Rate=k[Aldehyde]^2[\overset{\ominus}{O}H]^2}$ in concentrated base}
\opt{B}{Self cannizaro is a disproportionation reaction}
\opt{C}{Hydride transfer is the rate limiting step}
\opt{D}{$\mathrm{Ph-CHO}$ and $\mathrm{H-\overset{O}{\overset{\|}{C}}-H}$ gives $\mathrm{CH_3OH}$ + \qfigi[0.8]{p1_chem_s2_q11_1}}
\end{qoptions}
\qn{12} Which of the following are correct
\begin{qoptions}
\opt{A}{\qfigi{p1_chem_s2_q12_A}}
\opt{B}{\qfigi{p1_chem_s2_q12_B} can form stable hydrate}
\opt{C}{\qfigi{p1_chem_s2_q12_C} undergo cannizzaro reaction}
\opt{D}{\qfigi{p1_chem_s2_q12_D} mainly undergo self aldol condensation}
\end{qoptions}
\sect{SECTION-II}
\qn{13} An aqueous solution contains 10\% ammonia by mass and has a density of 0.85 g/ml. If $\left[\mathrm{H_3O^+}\right]$ in this solution is `x M', then the value of `$x\times10^{12}$' is ($K_a$ for $\mathrm{NH_4^+}=5.0\times10^{-10}$ M)
\qn{14} A solution of weak acid was titrated with base NaOH. The equivalence point was reached when 36.12 ml of 0.1 M-NaOH have been added. Now, 18.06 ml of 0.1 M-HCl was added to the titrated solution, the pH was found to be 5.0. The $\mathrm{p}K_a$ of acid is
\qn{15} $\mathrm{\underset{(A)}{C_8H_8O}\xrightarrow[\Delta]{dilOH^-}S\xrightarrow[H_3O^+]{CH_3MgBr\,(1\ mole)}(X)}$

A is aromatic carbonyl compound, which can give haloform reaction. Number of stereoisomer in the final product X.
\qfullbegin
\qn{16} \hfill
\qfig{p1_chem_s3_q16_A}
\qfullsplit
\qfig{p1_chem_s3_q16_2}

$X$ = Number of compound obtained by aldol reaction

$Y$ = Number of compounds react with $\mathrm{NaHCO_3}$

Sum of $X+Y$ is
\qfullend
\qn{17} How many of the following reactions are correct consider the following products are major

(i) $\mathrm{CH_3-\overset{O}{\overset{\|}{C}}-CH_3}$ (1 mole) $+$ Methanal ($x$ moles) $\longrightarrow$ \qfigi[0.8]{p1_chem_s3_q17_1} (1 mole). In the above reaction $x=6$

(ii) \qfigi[0.8]{p1_chem_s3_q17_2} $\xrightarrow[2)\,\mathrm{H_2O}]{1)\,\mathrm{LiAlH_4}}$ \qfigi[0.8]{p1_chem_s3_q17_3}

(iii) \qfigi[0.8]{p1_chem_s3_q17_4} $\xrightarrow[(2)\,\mathrm{H_2O}]{(1)\,\mathrm{NaBH_4}}$ $\underset{\text{(alcohol)}}{X}$

`$X$' can show optical isomerism, but does not give haloform reaction

(iv) $\mathrm{Ph-\underset{(P)}{\overset{O}{\overset{\|}{C}}}-C_2H_5}$ ;

`$P$' does not react with $\mathrm{NaHSO_3}$ but gives Iodoform reaction.
\qn{18} $\underset{(x\ \text{moles})}{\text{Methanal}}+\underset{(y\ \text{moles})}{\mathrm{NH_3}}\xrightarrow{\Delta}P$ (urinary antiseptic used as)

Then value of $\dfrac{2x}{y}$ is
\partsec{3}{MATHEMATICS}
\sect{SECTION-I (i)}
\qn{1} $\displaystyle\lim_{n\to\infty}\sum_{r=0}^{n}\left({}^nC_r\cdot\dfrac{1}{(r+4)\cdot n^r}\right)$ is equal to
\begin{qoptions}
\opt{A}{$2e+6$}
\opt{B}{$6-2e$}
\opt{C}{$6-e$}
\opt{D}{$6+e$}
\end{qoptions}
\qn{2} If $A(z_1),B(z_2),C(z_3)$ are the vertices of a triangle $ABC$ inscribed in the circle $|z|=2$. Internal angle bisector of the angle $A$, meet the circumcircle again at $D(z_4)$, then :
\begin{qoptions}
\opt{A}{$z_4^2=z_2z_3$}
\opt{B}{$z_4=\dfrac{z_2z_3}{z_1}$}
\opt{C}{$z_4=\dfrac{z_1z_2}{z_3}$}
\opt{D}{$z_4=\dfrac{z_1z_3}{z_2}$}
\end{qoptions}
\qn{3} Let $I(n)=\displaystyle\int_1^e x^3\cdot(\ell nx)^n\,dx$, where $n\in N$. Find the value of $\ell n\left(4I(5)+5I(4)\right)$.
\begin{qoptions}
\opt{A}{3}
\opt{B}{4}
\opt{C}{5}
\opt{D}{6}
\end{qoptions}
\qn{4} Let $f(x)=x^4+|x|$ and $I_1=\displaystyle\int_0^{\pi}f(\cos x)\,dx$ and $I_2=\displaystyle\int_0^{\pi/2}f(\sin x)\,dx$, then the value of $\dfrac{I_1}{I_2}$ is :
\begin{qoptions}
\opt{A}{1}
\opt{B}{$\dfrac12$}
\opt{C}{2}
\opt{D}{$\dfrac14$}
\end{qoptions}
\qn{5} The value of definite integral $\displaystyle\int_{1/3}^{2/3}\dfrac{\ln x}{\ln(x-x^2)}\,dx$ is equal to :
\begin{qoptions}
\opt{A}{$\dfrac13$}
\opt{B}{$\dfrac16$}
\opt{C}{$\dfrac19$}
\opt{D}{$\dfrac1{18}$}
\end{qoptions}
\qn{6} The evaluation of $\displaystyle\int\dfrac{px^{p+2q-1}-qx^{q-1}}{x^{2p+2q}+2x^{p+q}+1}\,dx$ is :
\begin{qoptions}
\opt{A}{$-\dfrac{x^p}{x^{p+q}+1}+C$}
\opt{B}{$\dfrac{x^q}{x^{p+q}+1}+C$}
\opt{C}{$-\dfrac{x^q}{x^{p+q}+1}+C$}
\opt{D}{$\dfrac{x^p}{x^{p+q}+1}+C$}
\end{qoptions}
\sect{SECTION-I (ii)}
\qn{7} For the function $f(x)=\displaystyle\int_0^{x^2}\left(\cos\dfrac1{\sqrt t}+1\right)dt,\ x\ge1$
\begin{qoptions}
\opt{A}{$f(x)$ is an increasing function in the interval $[1,\infty)$}
\opt{B}{$\displaystyle\lim_{x\to\infty}f''(x)=4$}
\opt{C}{$f(4)+f(6)>2f(5)$}
\opt{D}{For at least one $x$ in $[1,\infty)$, $f'(x+3)-f'(x)>12$}
\end{qoptions}
\qn{8} $\displaystyle\int_0^{\pi/2}\tan^3x\cos^5x\cos7x\,dx$ is equal to
\begin{qoptions}
\opt{A}{$\dfrac18\displaystyle\int_\pi^0\sin^2x\sin3x\,dx$}
\opt{B}{$\dfrac18\displaystyle\int_{\pi/2}^0\sin^2x\sin3x\,dx$}
\opt{C}{$\dfrac1{30}$}
\opt{D}{$\dfrac1{60}$}
\end{qoptions}
\qn{9} $S_n=\displaystyle\sum_{r=1}^{n}\left(\dfrac{r^4+r^3n+r^2n^2+2n^4}{n^5}\right)$ and $T_n=\displaystyle\sum_{r=0}^{n-1}\left(\dfrac{r^4+r^3n+r^2n^2+2n^4}{n^5}\right),\ n=1,2,3,\ldots$ then
\begin{qoptions}
\opt{A}{$T_n>\dfrac{167}{60}$}
\opt{B}{$T_n<\dfrac{167}{60}$}
\opt{C}{$S_n>\dfrac{167}{60}$}
\opt{D}{$S_n<\dfrac{167}{60}$}
\end{qoptions}
\qn{10} Let $f(x)$ be a function given by $f(x)=x^{3/2}e^{-1/2}$. Then --
\begin{qoptions}
\opt{A}{$f$ is non-monotonic in $(4,\infty)$}
\opt{B}{$\dfrac{9\sqrt3}{e\sqrt e}<\displaystyle\int_0^3f(x)\,dx<\dfrac{9\sqrt3}{\sqrt e}$}
\opt{C}{$\displaystyle\int_0^3f(x)\,dx<\dfrac{9\sqrt3}{e\sqrt e}$}
\opt{D}{$f$ is monotonic in $(4,\infty)$}
\end{qoptions}
\qn{11} Consider a real valued function $H(x)=\displaystyle\int_{-1}^{-1-x^2}(z+2)(z+5)(z+10)\,dz$, then
\begin{qoptions}
\opt{A}{$H(x)$ has 7 critical points}
\opt{B}{$H''(x)=0$ has exactly 3 positive roots}
\opt{C}{$H'''(x)=0$ has exactly 2 positive roots}
\opt{D}{$H'(x)H''(x)=0$ has 13 distinct real roots}
\end{qoptions}
\qn{12} If $z_1,z_2,z_3$ are any three roots of the equation $z^6=(z+1)^6$, then $\arg\left(\dfrac{z_1-z_2}{z_2-z_3}\right)$ can be equal to
\begin{qoptions}
\opt{A}{0}
\opt{B}{$\pi$}
\opt{C}{$\dfrac\pi4$}
\opt{D}{$-\dfrac\pi4$}
\end{qoptions}
\sect{SECTION-II}
\qn{13} Let $f(x)$ is a polynomial such that $f(x)+f''(x)=x^3+7x$ then $\displaystyle\lim_{n\to\infty}\dfrac{\displaystyle\sum_{r=1}^{n}f(2r)}{1.2^2+2.3^2+\ldots\ldots+n(n+1)^2}$ is equal to
\qn{14} Tangent is drawn at the point $(x_i,y_i)$ on the curve $y=f(x)$, which intersects the $x$-axis at $(x_{i+1},0)$. Now, again a tangent is drawn at $(x_{i+1},y_{i+1})$ on the curve which intersects the $x$-axis at $(x_{i+2},0)$ and process is repeated $n$ times i.e. $i=1,2,3,\ldots\ldots,n$. If $x_1,x_2,x_3,\ldots\ldots,x_n$ form an arithmetic progression with common difference equal to $\log_2e$ and curve passes through $(0,2)$. Now if curve passes through the point $(-2,k)$, then the value of $k$ is
\qn{15} If $\displaystyle\int_a^b\left(7x^2-9\right)dx=0,\ (a,b>0,\ a\ne b)$ then the maximum possible value of $(3a+2b)$ is
\qn{16} Let $P=\left|z^5+\bar z^3+6z\right|-2\left|z^4+1\right|$ and $|z|=1$ and $Q$ is number of real roots of $5-x=\sqrt{x+\sqrt{x+\sqrt x}}$, where $z\in C,x\in R$ then value of $P_{\max}+P_{\min}+Q$.

(where $P_{\max}$ and $P_{\min}$ are minimum and maximum value of $P$ respectively)
\qn{17} Let $\omega$ be a non real complex number satisfying $\omega^{23}=1$. If $N=\displaystyle\sum_{k=1}^{22}\dfrac{1}{1+\omega^{8k}+\omega^{16k}}$ then sum of digits of $N$ is
\qn{18} The value of the definite integral $\displaystyle\int_1^2\left(\dfrac{e^x-e^{2/x}}{x}\right)dx$ is
\end{multicols}
\end{document}
